\section{Solutions for Chapter~\ref{ch:tensor_networks}}

\paragraph{Exercise~\ref{ex:MPS-schmidt}.}
\phantomsection
\label{sol:ex:MPS-schmidt}
The largest two Schmidt probabilities sum to $5/6$. Their
normalized truncation is
\[
 |\Psi_2\rangle=\sqrt{3/5}|00\rangle+\sqrt{2/5}|11\rangle.
\]
Its squared overlap with $\Psi$ is $5/6$, its phase-aligned
vector-norm error is $\sqrt{2-2\sqrt{5/6}}$, and its energy
is $-5/6$. The best rank-one energy is $-1/2$.
The proof by the projector onto the left Schmidt support
already optimizes over every subspace of the prescribed
dimension. It therefore includes arbitrary local basis
changes; it is not restricted to coordinate truncation in
the original product basis.



\paragraph{Exercise~\ref{ex:MPS-edge}.}
\phantomsection
\label{sol:ex:MPS-edge}
The only nonzero block of $A^a$ in block row $a$ and column $b$
is $W_{ab}$. In $A^{x_0}A^{x_1}$, the first tensor's column
must therefore be the second tensor's row $x_1$, leaving
$W_{x_0x_1}$ times the next block row. Induction produces
$W_{x_0x_1}\cdots W_{x_{L-2}x_{L-1}}$ followed by the last
row. The ordinary trace forces its last column back to $x_0$,
giving $W_{x_{L-1}x_0}$ and the desired cyclic fiber trace.
Any forbidden edge makes one factor zero. This proof works
for unequal fiber dimensions, because consecutive rectangular
maps have matching intermediate spaces.



\paragraph{Exercise~\ref{ex:MPS-correlation}.}
\phantomsection
\label{sol:ex:MPS-correlation}
The norm matrix is $(t^2-1)I+J$. Its eigenvalues are $t^2+2$
on the uniform vector and $t^2-1$ twice on the color-orthogonal
plane. The nontrivial length is
\[
 \xi=-\left[\log\left|\frac{t^2-1}{t^2+2}\right|\right]^{-1}.
\]
At $t=0$ the ratio is $-1/2$, giving alternating color
correlations and finite length $1/\log2$. At $t=1$ it is zero;
the uniform edge amplitudes give an ordinary product state on
the supported three colors, with no nonzero-distance connected tail.
At large $|t|$,
$|\lambda|=1-3/(t^2+2)$ and $-\log|\lambda|\sim3/t^2$,
so $\xi\sim t^2/3$ in lattice sites. This growth reflects long
domains of equal colors. No conformal limit follows merely from this parameter
limit.



\paragraph{Exercise~\ref{ex:MPS-cat}.}
\phantomsection
\label{sol:ex:MPS-cat}
Across any nontrivial cut the three distinct product branches
are orthogonal on both sides. Their Schmidt probabilities are
$(1/3,1/3,1/3)$, so the entropy is $\log3$. It is independent
of interval length. The state has three ordered product components
and a nondecaying order correlator; its standard symmetry-breaking
parent has a finite ground-state degeneracy. Neither
feature gives the logarithmic interval growth associated with
a CFT central charge.



\paragraph{Exercise~\ref{ex:MPS-resolvent}.}
\phantomsection
\label{sol:ex:MPS-resolvent}
Normalize $(l|r)=1$ and put $\mathcal P=|r)(l|$. The fixed-point
relations imply $\mathcal E\mathcal P=\mathcal P\mathcal E=
\mathcal P$. Thus
$\mathcal E^n=\mathcal P+(\mathcal E-\mathcal P)^n$ for
$n\geq1$. Subtracting the product of one-point functions in
a separated correlator removes the $\mathcal P$ term.
On the complementary space the spectral radius is less than
one for an injective clustering tensor, so the geometric
series converges even if the map is not diagonalizable:
\[
 \sum_{n\geq0}(\mathcal E-\mathcal P)^n
 =(1-\mathcal E+\mathcal P)^{-1}.
\]
The additional stationary direction of a cat remains after
subtracting only one rank-one projector. A periodic peripheral
root likewise has modulus one. The series then fails to
converge on its coupled direction, and a connected primitive
component must be specified before using this formula.



\paragraph{Exercise~\ref{ex:MPS-scaling}.}
\phantomsection
\label{sol:ex:MPS-scaling}
At $n=1,2,\infty$, the logarithmic coefficients in
\eqref{eq:FES} are respectively $c/6,c/8,c/12$.
Let $x_a=\log\xi_a$ and $x_b=\log\xi_b$. For
$S_1=ax+b+u/x$, the secant is
\[
 \frac{S_1(x_b)-S_1(x_a)}{x_b-x_a}
 =a-\frac{u}{x_ax_b}.
\]
The corresponding central-charge diagnostic is six times this
quantity. Increasing $\xi$ by a large factor may still change
the logarithms modestly, so the correction can drift slowly.
This illustrative correction is not a derived law for the
Haagerup sequence; it shows why apparent slope stability needs
an independently justified remainder model.


